\documentclass[11pt,a4paper]{article}
\usepackage[T1]{fontenc}
\usepackage{lmodern,microtype}
\usepackage[margin=25mm]{geometry}
\usepackage{amsmath,amssymb,amsthm,mathrsfs,booktabs,array}
\usepackage[hidelinks]{hyperref}
\ifdefined\pdfinfoomitdate\pdfinfoomitdate=1\fi
\ifdefined\pdftrailerid\pdftrailerid{}\fi
\ifdefined\pdfsuppressptexinfo\pdfsuppressptexinfo=15\fi
\hypersetup{pdftitle={Sun's supercongruence for two Apery-like binomial sums at discriminant -20},pdfauthor={Oleksiy Babanskyy}}
\pagestyle{plain}
\newtheorem{theorem}{Theorem}[section]
\newtheorem{lemma}[theorem]{Lemma}
\newtheorem{proposition}[theorem]{Proposition}
\newtheorem{corollary}[theorem]{Corollary}
\theoremstyle{remark}
\newtheorem{remark}[theorem]{Remark}
\newcommand{\Z}{\mathbb Z}
\newcommand{\Q}{\mathbb Q}
\newcommand{\C}{\mathbb C}
\newcommand{\Fp}{\mathbb F_p}
\newcommand{\Fpp}{\mathbb F_{p^2}}
\newcommand{\Zp}{\mathbb Z_p}
\newcommand{\dd}{\partial}
\newcommand{\leg}[2]{\left(\frac{#1}{#2}\right)}
\newcommand{\smat}[4]{\begin{pmatrix}#1&#2\\#3&#4\end{pmatrix}}
\newcommand{\Gbase}{\Gamma_0(6)}
\DeclareMathOperator{\ord}{ord}
\DeclareMathOperator{\Res}{Res}
\DeclareMathOperator{\Disc}{Disc}
\setlength{\emergencystretch}{2em}
\allowdisplaybreaks
\title{Sun's supercongruence for two Ap\'ery-like binomial sums at discriminant $-20$}
\author{Oleksiy Babanskyy}
\date{}

\begin{document}
\maketitle
\begin{abstract}
Let $a_n=\sum_{j=0}^n\binom nj^2\binom{2j}j$ and let
$Q_n=\sum_{j=0}^n\binom nj(-8)^{n-j}\sum_{k=0}^j\binom jk^3$.
We prove Sun's Conjecture~3.7: for every prime $p\ne2,5$, the two sums
$(-1)^{(p-1)/2}\sum_{n<p}\binom{2n}n a_n20^{-n}$ and
$(-1)^{(p-1)/2}\sum_{n<p}\binom{2n}n Q_n(-16)^{-n}$ are congruent
modulo $p^2$ to $4x^2-2p$ if $p=x^2+5y^2$, to $2x^2-2p$ if
$2p=x^2+5y^2$, and to zero if $(-5/p)=-1$.
For split primes, two modular periods at one complex multiplication
point yield separate finite specializations of an integral Hecke
polynomial. Its degree-$p$ endpoint and the sign in the nonprincipal
class are treated explicitly. For inert primes, a degree-five
endomorphism with square $[-5]$ forces a Hasse polynomial to vanish;
a finite quadratic differential invariant then lifts each original sum
from modulo $p$ to modulo $p^2$. Fixed algebraic certificates and their
uniform specialization are included.
\end{abstract}

\section{The supercongruence and its two proof branches}\label{sec:intro}
For $n\ge0$ define the Franel numbers and two Ap\'ery-like sequences by
\begin{equation}\label{eq:sequences}
\begin{split}
 f_n&=\sum_{j=0}^n\binom nj^3,\qquad
 a_n=\sum_{j=0}^n\binom nj^2\binom{2j}j,\\
 Q_n&=\sum_{j=0}^n\binom nj(-8)^{n-j}f_j.
\end{split}
\end{equation}
Write
\begin{equation}\label{eq:series}
 G_a(t)=\sum_{n\ge0}\binom{2n}n a_nt^n,\qquad
 G_Q(t)=\sum_{n\ge0}\binom{2n}n Q_nt^n.
\end{equation}
Their truncations at degrees $0,\ldots,p-1$ are denoted by
$G_{a,p}(t)$ and $G_{Q,p}(t)$. This convention for the subscript $p$
will also be used for a generic series $G$. All congruences of
$p$-integral rational numbers are interpreted in $\Zp$.

\begin{theorem}\label{thm:main}
For every prime $p\ne2,5$,
\begin{equation}\label{eq:main}
\begin{aligned}
 &(-1)^{(p-1)/2}G_{a,p}(1/20)
 \equiv(-1)^{(p-1)/2}G_{Q,p}(-1/16)\\
 &\qquad\equiv
 \begin{cases}
 4x^2-2p,&p=x^2+5y^2,\quad p\equiv1,9\pmod{20},\\
 2x^2-2p,&2p=x^2+5y^2,\quad p\equiv3,7\pmod{20},\\
 0,&p\equiv11,13,17,19\pmod{20}
 \end{cases}
 \pmod{p^2}.
\end{aligned}
\end{equation}
\end{theorem}

This is Conjecture~3.7 of Zhi-Hong Sun
\cite[p.~2742]{sun}. The existence of the representations in the two
split cases is proved in Lemma~\ref{lem:forms}. The signs of $x,y$ do
not affect the displayed values.

The two sums are evaluated separately. For split primes we identify
$G_a,G_Q$ as modular periods, derive a finite polynomial congruence,
and evaluate it at a common point of discriminant $-20$. Both ideal
classes occur, giving matrices of determinant one and two. For inert
primes we first reduce the sums modulo $p$ to a square of a Legendre
polynomial. An explicit endomorphism of an elliptic curve forces that
polynomial to vanish, and a differential invariant lifts the vanishing
of each sum modulo $p^2$. Thus the equality of the two sums follows
from their evaluations; it is not used as a bridge in either argument.
The prime $3$ is computed directly at the end of Section~\ref{sec:inert}.

The finite invariant and integral Hecke specialization follow the
organization of the author's treatments of Sun's Conjectures~3.5 and~3.6
\cite{babanskyy35,babanskyy36}. Their needed forms are established here;
the two previous supercongruences are not assumptions of the proof.

The modular orbit argument uses Beukers's Lemma~3.7
\cite{beukers}. The degree-$p$ endpoint follows the coefficient
comparison in the proof of his Theorem~6.1, whose stronger congruence
has different hypotheses; the
modulo-$p^2$ statement required here is proved in
Theorem~\ref{thm:specialization}. The modulo-$p$ Legendre and Franel
antecedents are identified in Section~\ref{sec:modp}; the classical
Legendre square identity is derived here from its symmetric-square
equation. The endomorphism is certified by polynomial identities before
reduction. The appendices give exact arithmetic inputs and the bounds
which make those finite computations sufficient.

\section{A finite differential invariant and lifting}\label{sec:invariant}
The three integral sequences $a_n,Q_n,f_n$ satisfy
\cite[pp.~2729--2730]{sun}
\begin{equation}\label{eq:recurrence}
 (n+1)^2u_{n+1}=(\alpha n(n+1)+\beta)u_n-\gamma n^2u_{n-1},
 \qquad u_0=1,\quad u_1=\beta,
\end{equation}
with parameters
\begin{equation}\label{eq:parameters}
\begin{array}{c|rrr}
u_n&\alpha&\beta&\gamma\\\hline
a_n&10&3&9\\
Q_n&-17&-6&72\\
f_n&7&2&-8
\end{array}
\end{equation}
For any such sequence put $v_n=\binom{2n}n u_n$,
$V(t)=\sum_{n\ge0}v_nt^n$, $U(t)=V_p(t)$, and $\theta=t\,d/dt$.
The central binomial factor converts \eqref{eq:recurrence} into $LV=0$,
where
\begin{equation}\label{eq:operator}
 L=\theta^3-2t(2\theta+1)(\alpha\theta^2+\alpha\theta+\beta)
       +4\gamma t^2(\theta+1)(2\theta+1)(2\theta+3).
\end{equation}

\begin{lemma}\label{lem:invariant}
Let $p>3$ be prime, let $\alpha,\beta,\gamma\in\Z$, and let $u_n$
be an integral sequence satisfying \eqref{eq:recurrence}. Define
\begin{equation}\label{eq:invariant}
\begin{split}
 D(t)&=1-4\alpha t+16\gamma t^2,\\
 I(Y)&=D\{2Y\theta^2Y-(\theta Y)^2\}
      +(\theta D)Y\theta Y+4t(-\beta+3\gamma t)Y^2.
\end{split}
\end{equation}
Then $I(U)\equiv0\pmod{p^2}$ as a polynomial identity.
\end{lemma}
\begin{proof}
Leibniz's rule gives the universal identity
\begin{equation}\label{eq:invariant-derivative}
 \theta I(Y)=2YLY.
\end{equation}
In characteristic zero this implies $I(V)=0$, since its constant term
is zero. The residual of the truncation is exactly
\begin{equation}\label{eq:truncation-residual}
 LU=-p^3v_pt^p+4\gamma p(2p-1)(2p+1)v_{p-1}t^{p+1}.
\end{equation}
The factor $\binom{2p-2}{p-1}$ shows that $p\mid v_{p-1}$.
Consequently $LU\equiv0$ and $\theta I(U)\equiv0\pmod{p^2}$.

The latter congruence alone does not determine coefficients at degrees
divisible by $p$. We have $\deg I(U)\le2p$ and constant term zero.
Every coefficient at a degree not divisible by $p$ is zero modulo
$p^2$, because that degree is a unit in $\Zp$. The two remaining
coefficients, calculated before reduction, are
\begin{equation}\label{eq:resonant}
 [t^p]I(U)=-2p^2v_p,\qquad
 [t^{2p}]I(U)=4\gamma(2p-1)(2p+1)v_{p-1}^2.
\end{equation}
For the first equality compare with $I(V)=0$: the first omitted term
$v_pt^p$ contributes $2p^2v_p$ at degree $p$. The second equality
comes from the highest-degree term of $U$. Both are divisible by
$p^2$, completing the coefficientwise proof.
\end{proof}

\begin{lemma}\label{lem:lifting}
Under the hypotheses of Lemma~\ref{lem:invariant}, let $t_0\in\Zp$
with $t_0D(t_0)\in\Zp^\times$. If $U(t_0)\equiv0\pmod p$, then
$U(t_0)\equiv0\pmod{p^2}$.
\end{lemma}
\begin{proof}
Evaluate the invariant modulo $p$. It gives
$\theta U(t_0)=0$, so the reduction of $U$ has a root of multiplicity
at least two at the reduction of $t_0$. The multiplicity is exactly
two. Indeed, $U(0)=1$ and $\deg U<p$, so a higher multiplicity would
satisfy $3\le m<p$. The coefficient of the ordinary third derivative
in $L$ is $t^3D(t)$. In the local expansion of $LU$ at $t_0$, the
term of degree $m-3$ then has nonzero coefficient proportional to
\[
 t_0^3D(t_0)m(m-1)(m-2).
\]
The lower derivatives cannot cancel it, contradicting
\eqref{eq:truncation-residual} modulo $p$. Therefore
$\theta^2U(t_0)$ is a unit. In $I(U)(t_0)$ modulo $p^2$, the terms
quadratic in $U(t_0)$ and $\theta U(t_0)$ vanish, leaving
\[
 2D(t_0)\theta^2U(t_0)U(t_0)\equiv0\pmod{p^2}.
\]
Its coefficient of $U(t_0)$ is a unit, proving the assertion.
\end{proof}

For the three parameters used below,
\begin{equation}\label{eq:D-values}
 D_a(1/20)=-\frac{16}{25},\qquad
 D_Q(-1/16)=\frac54,\qquad D_f(1/16)=-\frac54.
\end{equation}
All are units when $p>5$. Here an ordinary point means a nonsingular
point of the differential equation; it does not assert that an
associated elliptic curve is ordinary.

\section{Modular data and the two periods}\label{sec:periods}
\subsection{Eta quotients and cusp divisors}
Let $q=e^{2\pi i\tau}$, $\eta_d=\eta(d\tau)$, and
$\dd=q\,d/dq=(2\pi i)^{-1}d/d\tau$. Define
\begin{align}\label{eq:eta-data}
 h&=\frac{\eta_2\eta_6^5}{\eta_1^5\eta_3},&
 z&=\frac{\eta_1^6\eta_6}{\eta_2^3\eta_3^2},&
 J&=(\eta_1\eta_2\eta_3\eta_6)^2,\\
 U_\eta&=z^2,& B_\eta&=hz^2,& V_\eta&=h^2z^2,\notag\\
 R&=\frac{\eta_2^{12}\eta_3^{12}}{\eta_1^{12}\eta_6^{12}},&
 v&=\frac{\eta_2^{24}}{\eta_1^{24}},&
 S&=\left(\frac{\eta_2\eta_6}{\eta_1\eta_3}\right)^6.\notag
\end{align}
For the eta-quotient criteria and cusp-order formula we use
\cite[p.~2]{rouse-webb}. For the quotients listed in the following table
(which uses $z^2$, not $z$), the exponent sums
$\sum dr_d$ and $\sum(6/d)r_d$ are multiples of $24$; the weights are
even and $\prod d^{r_d}$ is a rational square. The character on
$\Gbase$ is therefore trivial. At the cusps
$0,1/2,1/3,\infty$ the orders are as follows:
\begin{center}\small
\begin{tabular}{@{}lrrrrcl@{}}\toprule
 & $r_1$&$r_2$&$r_3$&$r_6$&weight&cusp orders\\\midrule
$h$&$-5$&1&$-1$&5&0&$(-1,0,0,1)$\\
$R$&$-12$&12&12&$-12$&0&$(-1,1,1,-1)$\\
$v$&$-24$&24&0&0&0&$(-3,3,-1,1)$\\
$U_\eta$&12&$-6$&$-4$&2&2&$(2,0,0,0)$\\
$B_\eta$&7&$-5$&$-5$&7&2&$(1,0,0,1)$\\
$V_\eta$&2&$-4$&$-6$&12&2&$(0,0,0,2)$\\
$J$&2&2&2&2&4&$(1,1,1,1)$\\
$S$&$-6$&6&$-6$&6&0&$(-1,1,-1,1)$\\\bottomrule
\end{tabular}
\end{center}
The eta function has no zeros in the upper half-plane. Since $h$ has
one simple pole on the compact curve $X_0(6)$, it is a global
coordinate of degree one on that curve.

\begin{lemma}\label{lem:eta}
The following identities hold:
\begin{align}
 R&=h^{-1}+17+72h=\frac{(1+8h)(1+9h)}h,\label{eq:R}\\
 v(1+8h)&=h(1+9h)^3,&
 S&=\frac{h(1+9h)}{1+8h},\label{eq:vS}\\
 \Theta_a&:=U_\eta+18B_\eta+72V_\eta=\dd\log S,\label{eq:theta-a}\\
 \Theta_Q&:=U_\eta-72V_\eta=-\dd\log R.\label{eq:theta-Q}
\end{align}
The forms $\Theta_a,\Theta_Q$ are nonzero at every cusp.
\end{lemma}
\begin{proof}
The only poles of $R$ are simple and lie at $h=0,\infty$, so
$R=a/h+b+ch$. The exact expansions
\[
 h=q+5q^2+19q^3+O(q^4),\qquad
 R=q^{-1}+12+78q+O(q^2)
\]
give $a=1,b=17,c=72$.

The function $v(1+8h)-h(1+9h)^3$ has a pole of order at most four at
the cusp $0$, a pole of order at most one at $1/3$, and no other poles.
Its coefficients through $q^5$ vanish by the finite expansion in
Appendix~\ref{app:modular}, so its zero at infinity has order at least
six. Equality of the total orders of zeros and poles for a nonzero
meromorphic function on a compact curve proves the first identity in
\eqref{eq:vS}.

By \eqref{eq:R}, the remaining two cusps have $h$-values $-1/8,-1/9$.
At $1/3$ the pole of $v$ forces $1+8h=0$ in the preceding identity;
hence $h(1/3)=-1/8$ and $h(1/2)=-1/9$. The rational function asserted
for $S$ has the same divisor as $S$ and the same initial term $q$.
Their ratio is therefore one.

Logarithmic derivatives of modular units are modular forms of weight
two. They have no interior poles. At a cusp of width $w$, write the
local expansion as $Q^r(c_0+O(Q))$, with $c_0\ne0$. The transformed
logarithmic derivative has constant term $r/w$. Here $r=\pm1$ for
both $R,S$, so this constant is nonzero. The differences in
\eqref{eq:theta-a} and \eqref{eq:theta-Q} belong to
$M_2(\Gbase)$, and their coefficients $q^0,q^1,q^2$ vanish.
The index of $\Gbase$ in $\operatorname{SL}_2(\Z)$ is $12$;
the valence bound is $2$. Thus both differences vanish identically.
\end{proof}

\subsection{Identities of weight eight}
\begin{lemma}\label{lem:period}
Take either $\Theta=\Theta_a$ and $(\alpha,\beta,\gamma)=(10,3,9)$,
or $\Theta=\Theta_Q$ and $(\alpha,\beta,\gamma)=(-17,-6,72)$.
Set $W=\Theta\dd\log J-2\dd\Theta$. Then
\begin{align}
 W^2-\Theta^4+4\alpha J\Theta^2-16\gamma J^2&=0,\label{eq:weight8-one}\\
 2\Theta\dd^2\Theta-3(\dd\Theta)^2
          -4\beta J\Theta^2+12\gamma J^2&=0.\label{eq:weight8-two}
\end{align}
\end{lemma}
\begin{proof}
We first show that the residuals are holomorphic modular forms of
weight eight. For $M\in\Gbase$, put $j=c\tau+d$ and
$s_0=c/(2\pi i)$. The chain rule gives
\begin{align*}
 (\dd\Theta)(M\tau)&=j^4\dd\Theta+2s_0j^3\Theta,\\
 (\dd^2\Theta)(M\tau)&=j^6\dd^2\Theta
                    +6s_0j^5\dd\Theta+6s_0^2j^4\Theta,\\
 (\dd\log J)(M\tau)&=j^2\dd\log J+4s_0j.
\end{align*}
The anomalous terms cancel: $W$ has weight four, and
$2\Theta\dd^2\Theta-3(\dd\Theta)^2$ has weight eight.
For a cusp representative $\sigma$, let
$F_\sigma=\Theta|_2\sigma$ and $H_\sigma=J|_4\sigma$.
The same cancellation gives
\[
 W|_4\sigma=F_\sigma\dd\log H_\sigma-2\dd F_\sigma.
\]
The transformed second expression is
$2F_\sigma\dd^2F_\sigma-3(\dd F_\sigma)^2$.
The series $F_\sigma$ is holomorphic and
$H_\sigma=Q^r$ times a unit, so $\dd\log H_\sigma$ is regular.
Derivatives preserve nonnegative powers of the local parameter.
Thus the residuals are holomorphic at all cusps. There are no
interior poles, because $J$ does not vanish there.

For each of the four residuals the exact coefficients of
$q^0,\ldots,q^8$ are zero; Appendix~\ref{app:modular} specifies the
arithmetic calculation. The classical valence bound
\cite[Theorem~1.1(1)]{craig} is $8\cdot12/12=8$.
Order at least nine therefore proves the identities.
\end{proof}

\begin{proposition}\label{prop:parameters}
Define
\[
 P(h)=h(1+8h)(1+9h),\quad
 C(h)=1+18h+72h^2,\quad K(h)=1-72h^2.
\]
The parameters and modular periods are
\begin{align}
 t_a&=\frac{J}{\Theta_a^2}=\frac{P(h)}{C(h)^2},&
 t_Q&=\frac{J}{\Theta_Q^2}=\frac{P(h)}{K(h)^2},\label{eq:modular-parameters}\\
 G_a(t_a(q))&=\Theta_a(q),&
 G_Q(t_Q(q))&=\Theta_Q(q).\label{eq:period-identification}
\end{align}
Both parameters lie in $q+q^2\Z[[q]]$, and
$t_Q=t_a/(1-36t_a)$.
\end{proposition}
\begin{proof}
Comparison of eta exponents gives $J=h^2Rz^4$, and \eqref{eq:R}
then yields the rational formulas. All series are integral, the
constant terms of $\Theta_a,\Theta_Q$ are one, and $J=q+O(q^2)$.
The identity $C^2-K^2=36P$ gives the relation between the parameters.

For either row put $e=\dd\log t=W/\Theta$, a formal unit.
Equation~\eqref{eq:weight8-one} gives
$e^2=D(t)\Theta^2$. Substitute $\dd=e\theta$ into
\eqref{eq:weight8-two}, using
\[
 \frac{2\theta e}{e}=\frac{\theta D}{D}
                         +\frac{2\theta\Theta}{\Theta}.
\]
It becomes $I(\Theta)=0$. These divisions take place among formal
units. By \eqref{eq:invariant-derivative}, $L\Theta=0$.
In characteristic zero the operator $L$ determines the coefficient
of degree $n\ge1$ from the preceding ones by division by $n^3$.
It consequently has a unique formal solution with constant term one.
By \eqref{eq:recurrence}, $G_a,G_Q$ are those solutions.
\end{proof}

The rational relation between $t_a$ and $t_Q$ describes the parameters.
No transformation of their truncations modulo $p^2$ is inferred from it.

\section{Characters and pole orders}\label{sec:characters}
Use the determinant-normalized slash operator
\[
 (Y|_2M)(\tau)=\det(M)(c\tau+d)^{-2}Y(M\tau),
 \qquad
 W_2=\smat{2}{-1}{6}{-2},\quad W_6=\smat{0}{-1}{6}{0}.
\]
Let $\Gamma=\Gamma_0(6)+\langle2,6\rangle$ denote the full
Atkin--Lehner extension at level six; its projective quotient by
$\Gbase$ has order four \cite[\S3]{beukers}. Scalar matrices act
trivially under this slash convention.

\begin{lemma}\label{lem:characters}
The actions on $h$ are
\[
 \phi_2(h)=-\frac{1+8h}{8(1+9h)},\qquad
 \phi_6(h)=\frac1{72h}.
\]
Both $t_a,t_Q$ are $\Gamma$-invariant. The characters of the forms,
trivial on $\Gbase$, satisfy
\[
\begin{array}{c|rr}
 &W_2&W_6\\\hline
 \Theta_a&-1&-1\\
 \Theta_Q&-1&+1
\end{array}
\]
\end{lemma}
\begin{proof}
The matrices show that $W_2$ exchanges
$\infty\leftrightarrow1/3$ and $0\leftrightarrow1/2$,
while $W_6$ exchanges $\infty\leftrightarrow0$ and
$1/2\leftrightarrow1/3$. Since $h$ is a global coordinate, three
values from $0,\infty,-1/9,-1/8$ determine each displayed action.
Rational substitution gives
\[
 R(\phi_2(h))=R(h)^{-1},\qquad R(\phi_6(h))=R(h),
\]
\[
 S(\phi_2(h))S(h)=S(\phi_6(h))S(h)=1/64.
\]
The same substitution proves invariance of both parameters.
If $F(M\tau)=c_0F(\tau)^e$, differentiation gives
$(\dd\log F)|_2M=e\dd\log F$. Apply this to
\eqref{eq:theta-a} and \eqref{eq:theta-Q}.
\end{proof}

\begin{proposition}\label{prop:poles}
Each of $t_a,t_Q$ is a Hauptmodul for $\Gamma$. On the base curve it
has two double poles, and the corresponding $\Theta$ has a simple zero
at each of those points and no cusp zeros. A $\Gamma$-invariant
function with holomorphic numerators and at most $m$ factors
$\Theta^{-1}$ is a polynomial in $t$ of degree at most
$\lfloor m/2\rfloor$.
\end{proposition}
\begin{proof}
Each rational function in \eqref{eq:modular-parameters} has degree
four without cancellation. Invariance and the projective index four
give degree one on the quotient curve. The poles of $t_a$ are
$h=-1/6,-1/12$; those of $t_Q$ are $h=\pm1/\sqrt{72}$.
They are double and distinct from all cusps. The corresponding forms
are $C(h)z^2,K(h)z^2$, and $z$ has no interior zeros. The factors
$C,K$ have simple zeros. The base group has no elliptic stabilizers:
an order-two stabilizer is impossible modulo $3$, and an order-three
stabilizer is impossible modulo $2$. Thus these are ordinary local
coordinates on the base curve. Cusp nonvanishing was proved in
Lemma~\ref{lem:eta}.

On the quotient a meromorphic function whose only possible pole is at
the pole of a Hauptmodul is a polynomial in that Hauptmodul. On the
base cover its pole order is at most $m$, while a polynomial of degree
$d$ in $t$ has pole order $2d$. Hence $2d\le m$. This calculation is
on the base cover and requires no assignment of integral orders at
orbifold points of the quotient.
\end{proof}

\section{The integral Hecke polynomial and finite specialization}\label{sec:hecke}
\begin{lemma}\label{lem:endpoint}
For every prime $p>3$,
\[
 \binom{2p}p a_p\equiv6,\qquad
 \binom{2p}p Q_p\equiv-12\pmod{p^2}.
\]
\end{lemma}
The congruence $Q_p\equiv Q_1=-6\pmod{p^2}$ is known: it is the
$n=1$ consequence, after reduction modulo $p^2$, of Mao's stronger
congruence for $Q_{np}$ modulo $p^3$ \cite[Theorem~1.4]{mao}.
We give the alternative deduction from Straub below to make the
sequence convention explicit.
\begin{proof}
We have
\[
 \binom{2p}p=2\prod_{k=1}^{p-1}(1+p/k)\equiv2\pmod{p^2},
\]
since $\sum_{k=1}^{p-1}k^{-1}\equiv0\pmod p$. Each interior
term in $a_p$ contains $\binom pk^2$, so
$a_p\equiv1+\binom{2p}p\equiv3\pmod{p^2}$.

For $Q_p$ the interior terms have only one factor $\binom pk$ and
cannot be discarded. Straub's sequence
\cite[(10), Theorem~3.2]{straub} is
\[
 A_F(n)=\sum_{k=0}^n(-1)^k8^{n-k}\binom nk f_k=(-1)^nQ_n.
\]
His theorem gives
$A_F(p^rn)\equiv A_F(p^{r-1}n)\pmod{p^{2r}}$ for primes
$p\ge3$ and $n,r\ge1$. Taking $n=r=1$ gives
$Q_p\equiv-6\pmod{p^2}$. Multiplication by the central binomial
coefficient proves the two assertions. This import concerns one
coefficient and does not assume either sum in Theorem~\ref{thm:main}.
\end{proof}

\begin{lemma}\label{lem:hecke}
For either pair $(t,\Theta)$ above and any prime $p>3$, there is a
polynomial
\[
 \mathscr T_p(X,t)=\sum_{m=0}^{p+1}C_m(t)X^m\in\Z[X,t]
\]
with the exact identity
\begin{equation}\label{eq:hecke-polynomial}
\begin{split}
 \mathscr T_p(X,t(\tau))={}&
 \left(1-\frac{\Theta(p\tau)}{\Theta(\tau)}X\right)\\
 &\quad\cdot\prod_{k=0}^{p-1}
 \left(p^2-\frac{\Theta((\tau+k)/p)}{\Theta(\tau)}X\right).
\end{split}
\end{equation}
Its coefficients satisfy
\begin{equation}\label{eq:hecke-properties}
 \deg C_m\le\lfloor m/2\rfloor,\qquad
 \mathscr T_p\equiv C_pX^p+C_{p+1}X^{p+1}\pmod{p^2},
 \qquad C_{p+1}\equiv1\pmod p.
\end{equation}
\end{lemma}
\begin{proof}
Beukers's Lemma~3.7 \cite[pp.~17--18]{beukers} applies to weight two,
a character trivial on the base group and defined on its
Atkin--Lehner extension, and $p\nmid6$. These hypotheses are supplied
by Lemma~\ref{lem:characters}. With representatives
\[
 D_\infty=\smat{p}{0}{0}{1},\qquad
 D_k=\smat{1}{k}{0}{p}\quad(0\le k<p),
\]
the extended group permutes the ratios $(\Theta|_2D)/\Theta$.
The product in \eqref{eq:hecke-polynomial} is
$p^{2p}\prod_D(1-(\Theta|_2D)/(p\Theta)X)$, so its coefficients
are invariant.

The numerators are holomorphic in the interior and at cusps. For the
cusp assertion, a rational matrix of positive determinant takes a
rational cusp to a rational cusp. In local coordinates at the source
and target it becomes triangular and replaces a local parameter by
a positive power times a nonzero constant. It creates no negative
powers. Since $\Theta$ does not vanish at any cusp,
Proposition~\ref{prop:poles} gives polynomiality and the degree bound.

For integrality, form the product first in
$\Z[\zeta_p][[s_1]]$, where $s_1=q^{1/p}$. The denominators are
units with constant term one. Invariance under
$s_1\mapsto\zeta_ps_1$ eliminates exponents not divisible by $p$;
Galois invariance makes the remaining coefficients rational.
They are algebraic integers and thus lie in $\Z[[q]]$. Since
$t=q+q^2\Z[[q]]$, successive subtraction of its powers at the lowest
nonzero degree shows $C_m\in\Z[t]$.

For $m<p$, every contribution to the coefficient of $X^m$ contains
at least one constant factor $p^2$. Dividing by $p^2$ before the
integrality argument leaves the same symmetries and algebraic
integrality. The quotient therefore belongs to $\Z[t]$.
Since $p+1$ is even,
\[
 C_{p+1}=
 \frac{\Theta(q^p)\prod_{k=0}^{p-1}\Theta(\zeta_p^ks_1)}
      {\Theta(q)^{p+1}}.
\]
Reduce the numerator product modulo $1-\zeta_p$:
$\Theta(s_1)^p\equiv\Theta(q)$. Its coefficients are integers and
$(1-\zeta_p)\Z[\zeta_p]\cap\Z=p\Z$. Also
$\Theta(q^p)\equiv\Theta(q)^p\pmod p$. It follows that
$C_{p+1}\equiv1$ as a series, and by the same triangular argument
as a polynomial.
\end{proof}

\begin{theorem}\label{thm:specialization}
For every prime $p>3$, let $G=G_a$ or $G_Q$ and let $t$ be its
corresponding parameter. Then
\begin{equation}\label{eq:finite-specialization}
 C_p(t)+C_{p+1}(t)G_p(t)\equiv0\pmod{p^2}
\end{equation}
as a polynomial identity. If $t_0\in\Zp$ and $\xi$ is an exact
$p$-adic unit root of $\mathscr T_p(X,t_0)$, then
$G_p(t_0)\equiv\xi\pmod{p^2}$.
\end{theorem}
\begin{proof}
The endpoint comparison follows the coefficient comparison in the
proof of Beukers's Theorem~6.1 \cite[p.~24]{beukers}. His theorem has different
hypotheses and proves a congruence modulo $p^3$. The case needed
here modulo $p^2$, with pole-bound exponent $1/2$, is proved directly
as follows.

The formal quotient $\mathcal Q=\Theta(q)/\Theta(q^p)$ is a unit
root of \eqref{eq:hecke-polynomial}. Dividing by
$\mathcal Q^p$ in $\Zp[[t]]$ gives
$C_p+C_{p+1}\mathcal Q\equiv0$. Write $C_{p+1}=1+pE$ with
$E\in\Z[t]$. Then
\[
 \mathcal Q\equiv-(1-pE)C_p\pmod{p^2},\qquad
 \deg((1-pE)C_p)\le\frac{p+1}2+\frac{p-1}2=p.
\]
The degree bound is $p$, so its last coefficient must be handled.
The integral inverse $q=q(t)$ gives
$t^\sigma=t(q(t)^p)=t^p+O(t^{p+1})$. By
Proposition~\ref{prop:parameters},
\[
 \mathcal Q=\frac{G(t)}{G(t^\sigma)}
   =G_p(t)+\bigl([u^p]G(u)-[u]G(u)\bigr)t^p+O(t^{p+1}).
\]
Here $[u]G(u)$ is $6$ for $G_a$ and $-12$ for $G_Q$.
Lemma~\ref{lem:endpoint} eliminates precisely this extra coefficient
modulo $p^2$. All coefficients of the polynomial of degree at most
$p$ are now congruent to those of $G_p$. Multiplication by
$C_{p+1}$ proves \eqref{eq:finite-specialization}, including every
degree of its left side.

We now use only polynomials. From
$\mathscr T_p(\xi,t_0)=0$ and \eqref{eq:hecke-properties},
\[
 \xi^p\{C_p(t_0)+C_{p+1}(t_0)\xi\}\equiv0\pmod{p^2}.
\]
Divide by the unit $\xi^p$, subtract
\eqref{eq:finite-specialization} evaluated at $t_0$, and divide
by $C_{p+1}(t_0)\equiv1\pmod p$. This proves the specialization.
No infinite series has been evaluated at a unit parameter.
\end{proof}

\section{The common CM point and both split classes}\label{sec:split}
\begin{lemma}\label{lem:cm}
Let $\delta=i\sqrt5$ and $\tau_0=(1+\delta)/6$. Then
\[
 h(\tau_0)=\frac{\sqrt5-2+i(\sqrt5+2)}{36},\qquad
 t_a(\tau_0)=\frac1{20},\qquad t_Q(\tau_0)=-\frac1{16}.
\]
Neither $\Theta_a$ nor $\Theta_Q$ vanishes at $\tau_0$.
\end{lemma}
\begin{proof}
We have $-1/\tau_0=6\tau_0-2$ and
$-1/(2\tau_0)=3\tau_0-1$. The eta transformations
\cite[\S23.18, (23.18.5)--(23.18.7)]{dlmf}
\[
 \eta(\tau+1)=e^{\pi i/12}\eta(\tau),\qquad
 \eta(-1/\tau)=\sqrt{-i\tau}\,\eta(\tau)
\]
with their standard holomorphic square root give
\[
 \eta_6=e^{\pi i/6}\sqrt{-i\tau_0}\,\eta_1,
 \qquad
 \eta_3=e^{\pi i/12}\sqrt{-2i\tau_0}\,\eta_2.
\]
The number $-i\tau_0$ lies in the right half-plane, so
$\sqrt{-2i\tau_0}=\sqrt2\sqrt{-i\tau_0}$ with no extra sign.
Hence
\[
 h(\tau_0)=e^{3\pi i/4}\frac{(-i\tau_0)^2}{\sqrt2}
           =\frac{1-i}{2}\tau_0^2.
\]
Expansion gives the displayed value of $h$. Substitution into
\eqref{eq:modular-parameters} verifies $20P(h)=C(h)^2$,
$16P(h)=-K(h)^2$, and $C(h)K(h)\ne0$ at this point.
Since $z$ has no interior zeros, the forms are also nonzero there.
\end{proof}

\begin{lemma}\label{lem:forms}
For $p>5$ with $(-5/p)=1$, there are positive integers $x,y$ such
that either
\[
 p=x^2+5y^2,\quad p\equiv1,9\pmod{20},
\]
or
\[
 2p=x^2+5y^2,\quad p\equiv3,7\pmod{20}.
\]
In the first case $x,y$ have opposite parity; in the second both are
odd. In both cases $\gcd(x,y)=1$ and exactly one is divisible by $3$.
\end{lemma}
\begin{proof}
Choose $r^2\equiv-5\pmod p$. The primitive positive form
$[p,2r,(r^2+5)/p]$ has discriminant $-20$ and represents $p$.
Among its unimodular transforms choose a minimal leading coefficient
$a$, then normalize $|b|\le a$. Minimality gives $a\le c$, whence
$20=4ac-b^2\ge3a^2$ and $a\le2$. The only reduced forms are
$[1,0,5]$ and $[2,\pm2,3]$. In the latter case multiplication by
two gives $(2u\pm v)^2+5v^2=2p$.
Congruences modulo $4,5,8$ assign the stated residue classes and
parities. Neither coordinate is zero; their signs may be chosen
positive. A common square factor would divide $p$ or $2p$; the parity
conditions exclude a factor two, so the representation is primitive.
Modulo three, two nonzero coordinates would give
$x^2+5y^2\equiv0$, impossible for $p>3$; two zero coordinates would
give $9\mid p$ or $9\mid2p$.
Finally quadratic reciprocity gives $(-5/p)=1$ precisely in the
four classes $1,3,7,9\pmod{20}$. No uniqueness of the representation
is needed below.
\end{proof}

\begin{proposition}\label{prop:split}
Theorem~\ref{thm:main} holds for every split prime $p>5$.
\end{proposition}
\begin{proof}
First let $p=x^2+5y^2$. The preceding parity and divisibility
conditions give $\gcd(x+y,6y)=1$. Choose $a,b\in\Z$ with
$a(x+y)+6by=1$, and set
\[
 M_1=\smat{a}{b}{-6y}{x+y}\in\Gbase,\qquad
 k=b(x-y)-ay.
\]
The denominator at $\tau_0$ is $x-y\delta$. Rationalizing gives
\[
 M_1\tau_0=\frac{\tau_0+k}{p}.
\]
For either form its character on $M_1$ is one, so the transformation
law exhibits the exact root
\begin{equation}\label{eq:principal-root}
 \xi_1=p^2\frac{\Theta(\tau_0)}{\Theta(M_1\tau_0)}
       =\frac{p^2}{(x-y\delta)^2}=(x+y\delta)^2
\end{equation}
of the specialized polynomial \eqref{eq:hecke-polynomial}.
Periodicity makes the representative $k\bmod p$ legitimate.

Next let $2p=x^2+5y^2$, where $x,y$ are odd. Here
$\gcd(x+y,3y)=1$. Choose $a,b\in\Z$ with
$a(x+y)+3by=1$, and put
\[
 M_2=\smat{2a}{b}{-6y}{x+y},\qquad
 k=\frac{b(x-y)}2-ay\in\Z.
\]
Then $\det M_2=2$, $M_2\tau_0=(\tau_0+k)/p$, and
\[
 M_2W_2^{-1}=\smat{-2a-3b}{a+b}{3(y-x)}{x-2y}\in\Gbase.
\]
Both forms have character $-1$ on this class. In our normalization
$\Theta(M\tau)=\chi(M)(c\tau+d)^2\Theta(\tau)/\det M$.
Thus the exact root is
\begin{equation}\label{eq:nonprincipal-root}
 \xi_2=-\frac{2p^2}{(x-y\delta)^2}
      =-\frac{(x+y\delta)^2}{2}.
\end{equation}
This establishes the negative sign and the factor $1/2$ before any
reduction.

These roots are algebraic identities in $\Q(\sqrt{-5})$.
For a split prime choose its embedding in $\Q_p$ so that
$x+y\delta$ is a unit. The two factors $x\pm y\delta$ cannot both
vanish modulo $p$, by primitivity and $p>5$, and their product is
$p$ or $2p$. Hence the other factor has valuation exactly one.
The same embedding works for both forms, and
Theorem~\ref{thm:specialization} applies at $1/20,-1/16\in\Zp$.

In the principal case the trace of $\xi_1$ is $4x^2-2p$.
In the nonprincipal case the trace of $\xi_2$ is $2p-2x^2$.
The nonunit conjugate root has valuation two and disappears modulo
$p^2$, so either sum is congruent to the respective trace.
Finally $(-1)^{(p-1)/2}=1$ in the principal residue classes and $(-1)^{(p-1)/2}=-1$
in the nonprincipal classes. This proves the two signed values.
\end{proof}

\section{Finite transport, Legendre polynomials and Hasse}\label{sec:modp}
The Hasse--Legendre relation below is already contained in Sun's
\cite[Theorem~2.1(ii) and equation~(4.1)]{sun-legendre-ii}; the latter
also treats the parameter as a polynomial variable. His Theorem~4.2
gives the quartic reduction \eqref{eq:quartic-square} at $m=-1024$.
Sun also evaluates $F_{16,p}$ modulo $p$ in the principal split classes
\cite[Theorem~2.10]{sun-franel}. We retain local derivations of the
identities needed here.

\begin{lemma}\label{lem:transport}
Let $p$ be odd, $N=(p-1)/2$, and
$c,t,u_0,\ldots,u_{p-1}\in\Zp$. Suppose $D_0=1-4ct$ is a unit,
and define $b_n=\sum_{j=0}^n\binom njc^{n-j}u_j$. Then
\begin{equation}\label{eq:transport}
 \sum_{n<p}\binom{2n}n b_nt^n
 \equiv D_0^N\sum_{j<p}\binom{2j}j u_j(t/D_0)^j\pmod p.
\end{equation}
\end{lemma}
\begin{proof}
For $n<p$ we have
$\binom{2n}n\equiv(-4)^n\binom Nn\pmod p$. Substitute this and
reorder the finite sum. For $j\le N$, the coefficient of $u_j$ is
\[
 \binom Nj(-4t)^j\sum_{m=0}^{N-j}\binom{N-j}m(-4ct)^m
   =\binom Nj(-4t)^jD_0^{N-j}.
\]
For $j>N$ both coefficients vanish modulo $p$. All summands are
$p$-integral and $t/D_0$ is $p$-integral, so the
calculation gives \eqref{eq:transport}.
\end{proof}

Sun's binomial identity \cite[Lemma~3.1]{sun} is
$a_n=\sum_j\binom nj f_j$. For $Q_n$ the analogous transform is
its definition. With
\[
 F_{16,p}=\sum_{n=0}^{p-1}\binom{2n}n f_n16^{-n},
\]
the pairs $(c,t)=(1,1/20),(-8,-1/16)$ have $D_0=4/5,-1$.
Consequently
\begin{equation}\label{eq:transport-values}
 G_{a,p}(1/20)\equiv\leg5p F_{16,p},\qquad
 G_{Q,p}(-1/16)\equiv(-1)^{(p-1)/2} F_{16,p}\pmod p.
\end{equation}
These specializations also follow from \cite[Theorem~3.4]{sun}.
They are used only modulo $p$: even a modulo-$p^2$ vanishing of
$R_p$ does not pass through them without an additional argument.

\begin{lemma}\label{lem:dixon}
For $p>5$,
\begin{equation}\label{eq:dixon-reduction}
 F_{16,p}\equiv(-1)^{(p-1)/2} R_p\pmod p,\qquad
 R_p=\sum_{j=0}^{p-1}\frac{(4j)!}{(j!)^4}(-1024)^{-j}.
\end{equation}
\end{lemma}
\begin{proof}
Put $N=(p-1)/2$.
Dixon's form of the Franel numbers, recorded in
\cite[p.~2729]{sun}, is $f_n=\sum_k\binom nk^2\binom{2k}n$.
Group the sum by $k$. Modulo $p$, all terms with $n>N$ vanish;
when $n\le N$ we also have $k\le N$. Each group can be completed
through $n=2k<p$, since the added terms have central binomial
coefficient divisible by $p$. The exact identity for a completed
group is
\[
\begin{split}
 &\sum_{n=k}^{2k}\binom{2n}n\binom nk^2\binom{2k}n16^{-n}\\
 &\qquad=\frac{\binom{2k}k^2}{16^k}
       \sum_{j=0}^k\frac{(-k)_j(k+1/2)_j}{(j!)^2}(-1/4)^j.
\end{split}
\]
At $n=k+r$, the ratio of successive summands is
$(k-r)(2k+2r+1)/(8(r+1)^2)$; the initial term is
$\binom{2k}k^2/16^k$. This proves the identity without evaluating
an infinite hypergeometric series.

Modulo $p$, $k+1/2\equiv k-N$ and
$\binom{2k}k^2/16^k\equiv\binom Nk^2$. On interchanging the sums,
the inner sum is the integer identity
\begin{equation}\label{eq:vandermonde}
 \sum_{k=0}^N\binom Nk^2\binom kj\binom{N-k}j
 =\binom Nj^2\binom{2N-2j}{N-2j},
\end{equation}
with right side zero if $2j>N$. Indeed the summand equals
$\binom Nj^2\binom{N-j}{k-j}\binom{N-j}k$, so Vandermonde's
identity applies. For $2j\le N$,
\[
 \binom{2N-2j}{N-2j}\equiv(-1)^{(p-1)/2}\binom N{2j}\pmod p,
\]
and
\[
 \binom Nj^2\binom N{2j}(-1/4)^j
 \equiv\frac{(1/4)_j(1/2)_j(3/4)_j}{(j!)^3}(-1/4)^j
 =\frac{(4j)!}{(j!)^4}(-1024)^{-j}.
\]
For $j>\lfloor p/4\rfloor$ the latter factorial quotient is
divisible by $p$, with $j!$ a unit. This proves
\eqref{eq:dixon-reduction}, including both boundaries of the sum.
\end{proof}

Let $P_n$ be the Legendre polynomial normalized by $P_n(1)=1$.
For $p>3$, $N=(p-1)/2$ and $d=\lfloor p/4\rfloor$, define
\begin{equation}\label{eq:hasse-definition}
 f_u(X)=X^3-\frac{3(3u+5)}2X+9u+7,
 \qquad\mathcal H_p(u)=[X^{p-1}]f_u(X)^N.
\end{equation}

\begin{lemma}\label{lem:hasse}
In $\Fp[u]$,
\begin{equation}\label{eq:hasse-legendre}
 \mathcal H_p(u)=\leg6p P_d(u).
\end{equation}
\end{lemma}
\begin{proof}
The exact translation
\[
 f_u(Z+2)=Z\left(Z^2+6Z+\frac92(1-u)\right)
\]
preserves the coefficient at degree $p-1$ modulo $p$.
Indeed the degree is at most $3N$, and the possible extra
contributions contain
$\binom{p-1+j}{p-1}\equiv0\pmod p$ for $1\le j\le N$.
Taking the coefficient of $Z^N$ in the quadratic raised to $N$
therefore gives
\[
 \mathcal H_p(u)=\sum_{j=0}^{d}\binom N{2j}\binom{2j}j
       6^{N-2j}(9/2)^j(1-u)^j.
\]
The number of quadratic and constant factors must both be $j$, which
explains the entire support of this sum. Modulo $p$,
\[
 \binom N{2j}\binom{2j}j
 \equiv4^j\frac{(1/4)_j(3/4)_j}{(j!)^2},\qquad
 6^N=\leg6p,\qquad 6^{-2j}(9/2)^j=8^{-j}.
\]
The unordered pair $\{-d,d+1\}$ reduces to $\{1/4,3/4\}$,
depending on $p\bmod4$. Rodrigues's formula gives
\[
 P_d(u)=\sum_{j=0}^d\frac{(-d)_j(d+1)_j}{(j!)^2}
                        \left(\frac{1-u}{2}\right)^j,
\]
which proves the polynomial identity. All factorial denominators are
units modulo $p$.
\end{proof}

\begin{lemma}\label{lem:legendre-square}
For every integer $n\ge0$, the following identity holds in $\Q[u]$:
\begin{equation}\label{eq:legendre-square}
 P_n(u)^2=\sum_{j=0}^n\binom nj\binom{n+j}j\binom{2j}j
                            \left(\frac{u^2-1}{4}\right)^j.
\end{equation}
\end{lemma}
\begin{proof}
Rodrigues's formula gives the Legendre equation
$(1-u^2)y''-2uy'+\lambda y=0$, where $\lambda=n(n+1)$,
as well as $P_n(-u)=(-1)^nP_n(u)$ and $P_n(1)=1$.
For $W=y^2$, differentiation and substitution give
\[
\begin{split}
 &(1-u^2)^2W'''-6u(1-u^2)W''\\
 &\qquad+\{6u^2-2+4\lambda(1-u^2)\}W'-4\lambda uW=0.
\end{split}
\]
Since $P_n^2$ is even, write it as $V(z)$ for $z=1-u^2$.
With $\vartheta=z\,d/dz$ the equation becomes
\[
 \{\vartheta^3-z(\vartheta+1/2)(\vartheta-n)
                                     (\vartheta+n+1)\}V=0.
\]
This change is verified in $\Q(u)$; the new left side is
$-(1-u^2)/(8u)$ times the preceding one. Clearing the intermediate
denominator gives a polynomial identity. With $V(0)=1$, its
coefficient recursion uniquely yields
\[
 [z^j]V=\frac{(-n)_j(n+1)_j(1/2)_j}{(j!)^3}.
\]
The coefficients terminate for $j>n$. Use
$(-n)_j(n+1)_j/(j!)^2=(-1)^j\binom nj\binom{n+j}j$ and
$(1/2)_j/j!=4^{-j}\binom{2j}j$ to obtain the result.
\end{proof}

\begin{corollary}\label{cor:quartic}
For $p>5$, $d=\lfloor p/4\rfloor$ and either root $s^2=5$ in
$\Fpp$,
\begin{equation}\label{eq:quartic-square}
 R_p=P_d(s/2)^2\quad\hbox{in }\Fpp.
\end{equation}
\end{corollary}
\begin{proof}
Set $n=d$ and $u=s/2$ in \eqref{eq:legendre-square} and reduce
modulo $p$. We have $1-u^2=-1/4$ and
$\{-d,d+1\}=\{1/4,3/4\}$ modulo $p$. The coefficient becomes
$(4j)!(-1024)^{-j}/(j!)^4$. Its terms for $d<j<p$ vanish modulo
$p$, giving exactly $R_p$.
\end{proof}

Both \eqref{eq:hasse-legendre} and \eqref{eq:legendre-square} are
polynomial identities before any value is substituted. They remain
valid in extension fields; no character sum over $\Fp$ has been
reinterpreted as a character sum over $\Fpp$.

\section{An explicit endomorphism and the inert branch}\label{sec:inert}
\subsection{The degree-five map}
Let $s^2=5$ and $\delta^2=-5$. These are distinct roles: $s$ enters
the coefficients of the curve, and $\delta$ enters the vertical
coordinate of the map. Initially work in
$\mathcal R=\Z[1/10,s]/(s^2-5)$, adjoining $\delta$ and its inverse
when defining the map. Put
\begin{equation}\label{eq:elliptic-curve}
 a=-\frac{30+9s}{4},\qquad b=\frac{14+9s}{2},\qquad
 f(X)=X^3+aX+b,\qquad E_s:\ Y^2=f(X).
\end{equation}
Its Hasse coefficient is $\mathcal H_p(s/2)$. Define
\begin{equation}\label{eq:kernel-polynomial}
 k(X)=X^2-\frac{5+3s}{2}X+\frac{65+42s}{20},
\end{equation}
\begin{equation}\label{eq:numerator}
\begin{split}
 n(X)={}&X^5-(5+3s)X^4+\frac{120+45s}{2}X^3
       -\frac{1055+405s}{4}X^2\\
       &\quad+\frac{7045+3588s}{16}X-\frac{1395+540s}{4},
\end{split}
\end{equation}
and $A=-n$, $B=5k^2$, $r=A/B$.

\begin{theorem}\label{thm:isogeny}
Over every field of characteristic $p>5$ containing $s,\delta$ as
above, the map
\begin{equation}\label{eq:isogeny}
 \Phi(X,Y)=\left(r(X),\frac{Yr'(X)}{\delta}\right),
 \qquad \Phi(O)=O
\end{equation}
defines a separable endomorphism of degree five of $E_s$ and
$\Phi^2=[-5]$.
\end{theorem}
\begin{proof}
We use four fixed polynomial identities in $\mathcal R[X]$.
Their explicit inputs and arithmetic verification are described in
Appendix~\ref{app:isogeny}. First, with $T=A'B-AB'$,
\begin{equation*}\label{eq:C0}
 fT^2+5(A^3B+aAB^3+bB^4)=0.\tag{C0}
\end{equation*}
Dividing by $-5B^4$ gives $f(X)r'(X)^2/\delta^2=f(r(X))$,
so the map lands on the same curve.

To determine the composition including the sign of the second
coordinate, use the classical division-polynomial normalization
\cite[\S\S5.5--5.6]{suth5}. The required polynomials are
\begin{align}\label{eq:division-polynomials}
 \psi_3&=3X^4+6aX^2+12bX-a^2,\notag\\
 F_4&=X^6+5aX^4+20bX^3-5a^2X^2-4abX-a^3-8b^2,\notag\\
 \psi_5&=32f^2F_4-\psi_3^3,\notag\\
 m&=X\psi_5^2-8fF_4\psi_3(\psi_5-4F_4^2),\\
 \psi_7&=\psi_5\psi_3^3-128f^2F_4^3,\notag\\
 \Omega_5&=4\psi_7F_4^2-\psi_3^3(\psi_5-4F_4^2)^2.\notag
\end{align}
Reducing the multiplication formula with $Y^2=f(X)$ gives
\begin{equation}\label{eq:multiplication-five}
 [5](X,Y)=\left(\frac m{\psi_5^2},
                         \frac{Y\Omega_5}{\psi_5^3}\right).
\end{equation}
Here $\psi_4=4YF_4$ and $\Omega_5=\omega_5/Y$ in that convention.
For a polynomial $F$ of degree at most five write
$F^{[5]}(U,V)=V^5F(U/V)$. The other three certificates are
\begin{align}
 A^{[5]}(A,B)&=m,\tag{C1}\label{eq:C1}\\
 B^{[5]}(A,B)&=\psi_5^2,\tag{C2}\label{eq:C2}\\
 m'\psi_5-2m\psi_5'&=5\Omega_5.\tag{C3}\label{eq:C3}
\end{align}
The first two give $r(r(X))=m/\psi_5^2$.
For the second coordinate the chain rule and (C3) give
\[
 (\Phi^2)_Y=\frac{Y}{\delta^2}(r\circ r)'
 =-\frac{Y}{5}\frac{m'\psi_5-2m\psi_5'}{\psi_5^3}
 =-\frac{Y\Omega_5}{\psi_5^3}.
\]
Thus both coordinates agree with $[-5]$ in
\eqref{eq:multiplication-five}.

It remains to check reduction and extension at all points. Exact
arithmetic gives
\begin{align}\label{eq:good-reduction}
 \Delta(E_s)&=729(s-2),&N(\Delta)&=-729^2,\notag\\
 \Res(n,k)&=\frac{3^{10}(25-11s)}{250},&
 N(\Res(n,k))&=\frac{3^{20}}{5^5},\\
 \Disc(k)&=\frac9{10}(5-s),&N(\Disc(k))&=\frac{81}{5}.\notag
\end{align}
All coefficient denominators involve only $2$ and $5$. For $p>5$
the discriminant and resultant are units for both choices of $s$:
their norms have numerator a power of three. Hence the curve is
smooth, $n,k$ stay coprime, and $r$ still has degree five. At infinity
$r(X)\sim-X/5$ and $r'(X)\sim-1/5$, so the rational map is
nonconstant and takes $O$ to $O$.

A rational map between smooth projective curves extends at its
omitted points, and a nonconstant morphism between elliptic curves
preserving the origin is a group homomorphism
\cite[Theorem~4.15 and Definition~4.20]{suth4}. The poles of $r$ are
therefore points mapped to $O$. From the diagram
$X\circ\Phi=r\circ X$, and $\deg X=2$, it follows that
$\deg\Phi=5$. Its differential is nonzero; more precisely
\[
 \Phi^*\left(\frac{dX}{2Y}\right)=\delta\frac{dX}{2Y}.
\]
Thus $\Phi$ is separable. Identities (C0)--(C3) reduce from
$\mathcal R$ at every $p>5$, proving the theorem uniformly.
\end{proof}

\subsection{Geometric torsion, Hasse vanishing and lifting}
\begin{proposition}\label{prop:inert}
For $p>5$ with $(-5/p)=-1$ and either root $s^2=5$ in $\Fpp$,
we have $\mathcal H_p(s/2)=0$. Consequently
$G_{a,p}(1/20)\equiv G_{Q,p}(-1/16)\equiv0\pmod{p^2}$.
\end{proposition}
\begin{proof}
Theorem~\ref{thm:isogeny} constructs $\Phi$ directly in
characteristic $p$, over a field containing $s$ and $\delta$.
If $E_s$ were ordinary, its group of geometric $p$-torsion points
would be cyclic of order $p$ \cite[p.~1]{suth13}. Every
endomorphism acts on that group by a scalar $c\in\Fp$, even when
the morphism is defined over an extension of $\Fp$. The relation
$\Phi^2=[-5]$ would force $c^2=-5$, impossible for an inert prime.
Thus $E_s$ is supersingular.

The Deuring--Hasse criterion \cite[p.~1]{debry} now gives
$[X^{p-1}]f(X)^{(p-1)/2}=0$, which is $\mathcal H_p(s/2)=0$.
If necessary one extends the finite field so that the smooth cubic
splits. Ordinarity, supersingularity and vanishing of this coefficient
are preserved under field extension. This argument applies to either
choice of $s$.

Lemma~\ref{lem:hasse}, Corollary~\ref{cor:quartic},
Lemma~\ref{lem:dixon}, and \eqref{eq:transport-values} give
\[
 \mathcal H_p(s/2)=0\ \Longrightarrow\ P_d(s/2)=0
 \ \Longrightarrow\ R_p=F_{16,p}=0\pmod p.
\]
Equation~\eqref{eq:transport-values} then gives
$G_{a,p}(1/20)\equiv G_{Q,p}(-1/16)\equiv0\pmod p$.
Apply Lemma~\ref{lem:lifting} separately to $G_{a,p}$ at $1/20$
and to $G_{Q,p}$ at $-1/16$. The unit conditions are exactly
\eqref{eq:D-values}. This yields vanishing modulo $p^2$ of each
original sum.
\end{proof}

\begin{proof}[Completion of Theorem~\ref{thm:main}]
Propositions~\ref{prop:split} and \ref{prop:inert} cover every
prime $p>5$. For $p=3$, the three summands in each of
$G_{a,3}(1/20)$ and $G_{Q,3}(-1/16)$ are $(1,3,0)$ modulo nine,
by their definitions. Each sum is therefore congruent to $4$;
multiplying by $(-1)^{(3-1)/2}=-1$ gives $5$ modulo nine.
The representation $2\cdot3=1^2+5\cdot1^2$ gives
$2\cdot1^2-2\cdot3\equiv5\pmod9$, as required.
The primes $2,5$ are excluded in the original statement.
\end{proof}

The supersingularity argument uses geometric $p$-torsion; it does not
assert that the integer Frobenius trace over $\Fpp$ is zero.
No complex uniformization or reduction theorem for a characteristic-zero
endomorphism enters this inert branch.

\section{Two consequences}\label{sec:consequences}
\begin{corollary}\label{cor:factor}
Let $p>5$ and $d=\lfloor p/4\rfloor$. In $\Fp[u]$,
\[
 4u^2-5\mid P_d(u)\quad\Longleftrightarrow\quad(-5/p)=-1.
\]
When the divisibility holds, both roots of $4u^2-5$ are simple roots
of $P_d$ over $\Fpp$.
\end{corollary}
\begin{proof}
For an inert prime, Proposition~\ref{prop:inert} and
Lemma~\ref{lem:hasse} give vanishing at both distinct roots
$u=\pm s/2$. At either point $1-u^2=-1/4$ is a unit.
In the Legendre equation, a multiple root of multiplicity
$2\le m\le d<p$ would produce an uncancelled leading local term
$m(m-1)(1-u^2)$. Thus the roots are simple.

Conversely, if $4u^2-5$ divides $P_d$, then
Corollary~\ref{cor:quartic}, Lemma~\ref{lem:dixon} and
\eqref{eq:transport-values} give
$G_{a,p}(1/20)\equiv0\pmod p$. At a split prime,
Theorem~\ref{thm:main} gives $(-1)^{(p-1)/2} G_{a,p}(1/20)\equiv4x^2$ or $2x^2$
modulo $p$, which is nonzero by the primitive representation.
Hence the prime is inert. This converse uses the already proved
theorem and is not a premise of its proof.
\end{proof}

\begin{corollary}\label{cor:stability}
If $p>5$ is inert, then for every $r\in\Zp$,
\[
 G_{a,p}(1/20+pr)\equiv0,\qquad
 G_{Q,p}(-1/16+pr)\equiv0\pmod{p^2}.
\]
\end{corollary}
\begin{proof}
The values at the unshifted parameters vanish modulo $p^2$ by
Proposition~\ref{prop:inert}. In Lemma~\ref{lem:lifting} their first
derivatives vanish modulo $p$, since the parameters are units and
$\theta U(t_0)\equiv0$. Taylor expansion of the finite polynomials
therefore proves the result.
\end{proof}

\appendix
\section{Finite modular certificates}\label{app:modular}
This appendix specifies the finite arithmetic underlying
Lemmas~\ref{lem:eta} and \ref{lem:period}. For an eta quotient
with exponent vector $(r_d)_{d\mid6}$, write
\[
 \prod_d\eta_d^{r_d}
 =q^{\sum dr_d/24}\prod_{m\ge1}(1-q^m)^{e_m},
 \qquad e_m=\sum_{d\mid m}r_d.
\]
To calculate through degree $M$, after removal of the initial power,
retain only $m\le M$ and use the finite binomial expansion
\[
 (1-q^m)^{e_m}
   =\sum_{j=0}^{\lfloor M/m\rfloor}(-1)^j\binom{e_m}j q^{mj}
                 \pmod{q^{M+1}}.
\]
For negative exponents the binomial coefficient has its usual integer
meaning $\binom{-e}j=(-1)^j\binom{e+j-1}j$. Multiply these finite
polynomials, discarding degrees exceeding $M$ after each step.

A second exact recurrence gives the same coefficients. If the
initial power has been removed and the constant term is one, set
\[
 \ell_n=-\sum_{m\mid n}m e_m,\qquad
 a_0=1,\qquad n a_n=\sum_{j=1}^n\ell_j a_{n-j}.
\]
This follows from the logarithmic derivative, so the divisions are
exact. Differentiation of a series uses $\dd q^n=nq^n$;
division by a series with unit constant term uses its triangular
coefficient recurrence. For $\dd\log J$, remove its initial $q$
and add the resulting constant one.

The following coefficients provide intermediate values:
\begin{center}\small
\begin{tabular}{@{}lrrrrr@{}}\toprule
 & $q^0$&$q^1$&$q^2$&$q^3$&$q^4$\\\midrule
$h$&0&1&5&19&61\\
$J$&0&1&$-2$&$-3$&4\\
$\Theta_a$&1&6&6&42&6\\
$\Theta_Q$&1&$-12$&$-12$&$-12$&$-12$\\
$t_a$&0&1&$-14$&117&$-884$\\
$t_Q$&0&1&22&405&6820\\
$v$&0&1&24&300&2624\\
$S$&0&1&6&21&68\\\bottomrule
\end{tabular}
\end{center}
The two form vectors through degree eight are
\[
 ([q^n]\Theta_a)_{n=0}^8=(1,6,6,42,6,36,42,48,6),
\]
\[
 ([q^n]\Theta_Q)_{n=0}^8=(1,-12,-12,-12,-12,-72,-12,-96,-12).
\]
Apply the preceding finite arithmetic to the exponent vectors in
Section~\ref{sec:periods}. For the named residuals the complete
required vectors are
\begin{center}\small
\begin{tabular}{@{}p{.66\textwidth}cl@{}}\toprule
Residual&degrees&vector\\\midrule
$v(1+8h)-h(1+9h)^3$&$0$--$5$&$(0,0,0,0,0,0)$\\
$\Theta_a-\dd\log S$&$0$--$2$&$(0,0,0)$\\
$\Theta_Q+\dd\log R$&$0$--$2$&$(0,0,0)$\\
Left side of \eqref{eq:weight8-one}, $\Theta_a$&$0$--$8$&$\mathbf0_9$\\
Left side of \eqref{eq:weight8-two}, $\Theta_a$&$0$--$8$&$\mathbf0_9$\\
Left side of \eqref{eq:weight8-one}, $\Theta_Q$&$0$--$8$&$\mathbf0_9$\\
Left side of \eqref{eq:weight8-two}, $\Theta_Q$&$0$--$8$&$\mathbf0_9$\\\bottomrule
\end{tabular}
\end{center}
Here $\mathbf0_9$ means nine zeros. These are $48$ coefficients:
six for a function with total pole order at most five, six for forms
of weight two with valence bound two, and $36$ for forms of weight
eight with valence bound eight. The three initial coefficients used
to determine $R$ were given in Lemma~\ref{lem:eta}.
Holomorphy and the pole bounds were proved before these coefficients
were used. Consequently each finite zero vector proves an identity,
rather than merely testing a prefix of an unbounded series.

\section{Fixed algebraic certificates}\label{app:isogeny}
The four identities (C0)--(C3) have only the fixed polynomials
\eqref{eq:elliptic-curve}--\eqref{eq:numerator} and
\eqref{eq:division-polynomials} as inputs. Their expansion can be
performed without algebraic-number approximations. Represent a
coefficient $u+vs$ by the pair $(u,v)\in\Q^2$, with rules
\[
 (u,v)+(u',v')=(u+u',v+v'),\qquad
 (u,v)(u',v')=(uu'+5vv',uv'+vu').
\]
Represent a polynomial by its coefficient list in ascending order.
Addition is componentwise, multiplication is finite convolution,
and differentiation sends its $i$th coefficient to $i$ times that
coefficient in degree $i-1$. For homogenization use the finite sum
\[
 F^{[5]}(A,B)=\sum_{i=0}^5 F_iA^iB^{5-i},
\]
padding $F$ with zeros to degree five. Thus every arithmetic operation
in the certificates is explicitly determined over rational pairs.

Form the residuals
\begin{align*}
 H_0&=f(A'B-AB')^2+5(A^3B+aAB^3+bB^4),\\
 H_1&=A^{[5]}(A,B)-m,\\
 H_2&=B^{[5]}(A,B)-\psi_5^2,\\
 H_3&=m'\psi_5-2m\psi_5'-5\Omega_5.
\end{align*}
Their degree bounds and full coefficient-pair vectors are
\begin{center}
\begin{tabular}{@{}lccc@{}}\toprule
Residual&degree at most&pairs&coefficient-pair vector\\\midrule
$H_0$&19&20&$((0,0))_{i=0}^{19}$\\
$H_1$&25&26&$((0,0))_{i=0}^{25}$\\
$H_2$&24&25&$((0,0))_{i=0}^{24}$\\
$H_3$&36&37&$((0,0))_{i=0}^{36}$\\\bottomrule
\end{tabular}
\end{center}
There are $108$ rational pairs, or $216$ rational entries.
The bounds follow directly from the input degrees: $\deg A=5$,
$\deg B=4$, $\deg f=3$, $\deg\psi_5=12$, $\deg m\le25$, and
$\deg\Omega_5\le36$. The pair arithmetic above gives the displayed
zero vectors. The companion supplies all inputs, intermediate
polynomials and exact expansion code; a separate symbolic expansion
provides a second implementation of these equalities.

For completeness, the nondegeneracy data in
\eqref{eq:good-reduction} follow from the standard polynomial discriminant
and resultant, with norm
$N(u+vs)=u^2-5v^2$. The short Weierstrass discriminant is
$-16(4a^3+27b^2)$, giving $729(s-2)$. These computations are again
over $\Q^2$. No unlisted prime denominator is introduced.

The coefficients of the fixed polynomials belong to $\mathcal R$;
their identities therefore hold after every specialization in
characteristic $p>5$. The unit resultant prevents cancellation of
the map's numerator and denominator after reduction, and the unit
curve discriminant ensures smoothness. This proves why the finite
certificates imply a statement for all the primes in
Theorem~\ref{thm:isogeny}. A finite panel of point counts or
supercongruence tests would not supply this implication.

\section{Companion and reproducibility}\label{app:companion}
The companion repository for this manuscript, version~0.1.1, is
\begin{center}
\url{https://github.com/aconsciousfractal/Sun-3-7-Supercongruence}.
\end{center}
It includes the standalone source, modular coefficient calculations,
the algebraic certificates (C0)--(C3), the endpoint calculations, and
finite checks of the original binomial definitions. Its deterministic
evidence manifest is identified by SHA-256
\begin{center}\ttfamily\small
\detokenize{035baf66ce55e008c441b05e316c54b323ad6b2146f5bc6f23eeede99d3ea977}
\end{center}
The evidence digest binds the stated arithmetic outputs and their
inputs; it does not certify novelty or mathematical correctness.
The modular identities and fixed algebraic certificates are used
with the sufficiency arguments in the preceding appendices.
Numerical checks at individual primes serve as regressions and are
not premises of the uniform theorem.

\paragraph{Acknowledgment.}
The author thanks Zhi-Hong Sun for comments on the presentation of
the companion manuscripts and for pointing out Mao's result on $Q_{np}$.

\paragraph{AI assistance and responsibility.}
AI assistance under the author's direction was used for mathematical
exploration, proof criticism, source comparison, manuscript drafting,
exact computation and repository preparation. OpenAI Codex and GPT-6.1 Sol
assisted with consolidation and separate checks of the proof and companion.
Research supplied through an external agent was also examined; its complete
model provenance is not asserted. The systems had access to related
arguments. Their checks are not independent human peer review, formal proof
verification or a determination of priority. Individual proof steps are not
attributed to particular systems. The author retains responsibility for
the manuscript, attributions and decision to publish.

\begin{thebibliography}{99}
\bibitem{sun}
Zhi-Hong Sun,
\emph{Supercongruences involving Ap\'ery-like numbers and binomial coefficients},
AIMS Mathematics \textbf{7}(2) (2022), 2729--2781.
\href{https://doi.org/10.3934/math.2022153}{doi:10.3934/math.2022153}.

\bibitem{sun-legendre-ii}
Zhi-Hong Sun,
\emph{Congruences concerning Legendre polynomials II},
Journal of Number Theory \textbf{133} (2013), 1950--1976.
\href{https://doi.org/10.1016/j.jnt.2012.11.004}{doi:10.1016/j.jnt.2012.11.004}.
Theorem and equation numbers refer to
\href{https://arxiv.org/abs/1012.3898v2}{arXiv:1012.3898v2}.

\bibitem{sun-franel}
Zhi-Hong Sun,
\emph{Congruences for sums involving Franel numbers},
International Journal of Number Theory \textbf{14} (2018), 123--142.
\href{https://doi.org/10.1142/S1793042118500094}{doi:10.1142/S1793042118500094}.

\bibitem{babanskyy35}
Oleksiy Babanskyy,
\emph{Sun's supercongruence for an Ap\'ery-like sequence}.
Manuscript and
\href{https://github.com/aconsciousfractal/Sun-3-5-Supercongruence}{companion repository}.

\bibitem{babanskyy36}
Oleksiy Babanskyy,
\emph{Sun's supercongruence for two Ap\'ery-like binomial sums}.
Manuscript and
\href{https://github.com/aconsciousfractal/Sun-3-6-Supercongruence}{companion repository}.

\bibitem{beukers}
Frits Beukers,
\emph{Supercongruences using modular forms},
\href{https://arxiv.org/abs/2403.03301v3}{arXiv:2403.03301v3}.
Lemma~3.7 and Theorem~6.1.

\bibitem{straub}
Armin Straub,
\emph{Gessel--Lucas congruences for sporadic sequences},
\href{https://arxiv.org/abs/2301.12248v1}{arXiv:2301.12248v1}.
Theorem~3.2.

\bibitem{rouse-webb}
Jeremy Rouse and John J. Webb,
\emph{On spaces of modular forms spanned by eta-quotients},
\href{https://arxiv.org/abs/1311.1460v3}{arXiv:1311.1460v3}.

\bibitem{mao}
Guo-Shuai Mao,
\emph{Congruences involving Franel numbers and Ap\'ery-like numbers},
preprint (2023), Theorem~1.4.
\href{https://doi.org/10.13140/RG.2.2.23378.73922}{doi:10.13140/RG.2.2.23378.73922}.
The cited author-uploaded full text is dated 8 January 2025 on ResearchGate.

\bibitem{craig}
William Craig,
\emph{New Types of Sturm bounds via p-adic transfer methods},
\href{https://arxiv.org/abs/2602.10240v1}{arXiv:2602.10240v1}.
Theorem~1.1(1), the classical valence bound.

\bibitem{dlmf}
NIST Digital Library of Mathematical Functions,
\href{https://dlmf.nist.gov/23.18}{Section~23.18},
transformations of the Dedekind eta function.

\bibitem{suth4}
Andrew V. Sutherland,
\emph{18.783 Elliptic Curves, Lecture~4}, MIT, Fall 2023.
\url{https://math.mit.edu/classes/18.783/2023/LectureNotes4.pdf}.

\bibitem{suth5}
Andrew V. Sutherland,
\emph{18.783 Elliptic Curves, Lecture~5}, MIT, Fall 2023.
\url{https://math.mit.edu/classes/18.783/2023/LectureNotes5.pdf}.

\bibitem{suth13}
Andrew V. Sutherland,
\emph{18.783 Elliptic Curves, Lecture~13}, MIT, Fall 2023.
\url{https://math.mit.edu/classes/18.783/2023/LectureNotes13.pdf}.

\bibitem{debry}
Christophe Debry,
\emph{Beyond two criteria for supersingularity: coefficients of division polynomials},
\href{https://arxiv.org/abs/1303.5002v1}{arXiv:1303.5002v1}.
The Deuring criterion.
\end{thebibliography}
\end{document}
